\documentclass{article}
\usepackage[utf8]{inputenc}
\usepackage{lmodern}
\usepackage[T1]{fontenc}
\usepackage{amsmath,amssymb}
\usepackage{graphicx}
\usepackage{hyperref}
\usepackage{booktabs}
\usepackage{natbib}
\usepackage[margin=1in]{geometry}

\title{Playing Pok\'{e}mon Red with a Language Model: An Autonomous LLM Agent for Retro RPG Gameplay}

\author{
  J\'{e}sica Collado \\
  \texttt{collado.jesica@gmail.com}
}

\date{}

\begin{document}

\maketitle

\begin{abstract}
We present an autonomous agent capable of playing Pok\'{e}mon Red, a classic role-playing game for the Game Boy, using a large language model as its sole decision-making engine. Unlike prior work that applies reinforcement learning or targets isolated game mechanics such as competitive battling, our agent plays the RPG through its opening act --- navigating an open world, engaging in combat, capturing Pok\'{e}mon, managing resources, and pursuing narrative objectives up to the S.S.\ Anne --- all through natural language reasoning over structured game state, with no task-specific training or fine-tuning.

The system is built around seven specialized LLM-based agents --- overworld navigation, battle strategy, catch-or-defeat evaluation, move learning, shop optimization, party lead selection, and nickname generation --- each operating on prompts constructed from direct memory reads of the emulator's working RAM. A minimal layer of mechanical automation handles low-level execution (menu traversal, pathfinding, dialogue advancement), while a feedback-based validation mechanism allows the LLM to self-correct when it attempts invalid actions. This separation between strategic reasoning and mechanical execution allows us to study which decisions an LLM can reliably make and how effectively it recovers from errors through feedback.

We evaluate the system across five open-weight models and one managed endpoint from the Qwen family, ranging from 9B to 397B parameters --- Qwen 3.5 9B, 27B, 35B-A3B, 122B-A10B, 397B-A17B, and Qwen 3.6 Plus --- all served via cloud inference through OpenRouter, to study how model scale affects gameplay quality. We discuss the design trade-offs involved in building a viable LLM-based game agent, the failure modes we observe, and the minimal set of interventions needed to keep the agent on track.
\end{abstract}

% --------------------------------------------------------------------------- %

\section{Introduction}
\label{sec:introduction}

Large language models have demonstrated remarkable capabilities as general-purpose reasoners, leading to growing interest in deploying them as autonomous agents in interactive environments. Recent work has shown that LLMs can navigate web interfaces, write and debug code, and operate software tools through natural language reasoning alone \citep{xi2023rise, wang2024survey}. A natural question arises: can these models handle the sustained, multi-faceted decision-making required to play a complete video game?

Video games, and role-playing games in particular, present a uniquely challenging testbed for language model agents. Unlike single-task benchmarks, an RPG demands the simultaneous coordination of multiple interleaved skills: spatial navigation through interconnected maps, turn-based combat with type-based strategy, resource management across items and currency, team composition and long-term planning, and the pursuit of narrative objectives conveyed through in-game dialogue. These decisions unfold over long time horizons, where early choices --- which Pok\'{e}mon to catch, when to heal, what items to stock --- compound into consequences hours later.

The dominant paradigm for game-playing AI has been reinforcement learning, which has achieved superhuman performance in domains ranging from Atari games \citep{mnih2015human} to Go \citep{silver2016mastering}. However, RL agents typically require millions of episodes of training, learn narrow policies tied to specific reward functions, and offer limited interpretability into their decision-making process. More recently, LLM-based agents have been applied to games such as Minecraft \citep{wang2023voyager} and competitive Pok\'{e}mon battles \citep{hu2024pokellmon}, but these efforts have either targeted open-ended sandbox environments or focused on isolated game mechanics rather than the full RPG experience.

We choose Pok\'{e}mon Red (1996) for the Game Boy as our testbed for several reasons. First, its mechanics are extensively documented by the fan community, including full maps of working RAM (WRAM) addresses, which enables precise programmatic access to game state without relying on computer vision. Second, the game offers a gradually increasing complexity curve: early routes are largely linear and forgiving, while later areas introduce branching paths, puzzle-based dungeons, and tougher opponents. Third, Pok\'{e}mon Red requires all the core RPG skills simultaneously --- exploration, combat, resource management, and quest tracking --- making it a comprehensive benchmark for evaluating an LLM's ability to coordinate multiple reasoning modes.

In this work, we present an autonomous agent that plays Pok\'{e}mon Red using a large language model as its sole decision-making engine, with no reinforcement learning, fine-tuning, or task-specific training. The system is organized around seven specialized LLM-based agents, each responsible for a distinct aspect of gameplay: overworld navigation, battle strategy, catch-or-defeat evaluation, move learning, shop optimization, party lead selection, and nickname generation. Each agent operates on structured prompts constructed from direct reads of the emulator's working RAM, and produces decisions in a constrained format that a mechanical execution layer translates into button presses.

A key design principle of our approach is the explicit separation between \emph{strategic reasoning} and \emph{mechanical execution}. The LLM is responsible for all high-level decisions --- where to go, what move to use, whether to catch a Pok\'{e}mon --- while a thin automation layer handles low-level tasks that require no intelligence: pathfinding via A*, menu traversal, and dialogue advancement. When the LLM attempts an invalid action, the system provides feedback explaining the failure and requests a new decision, preserving the LLM as the sole decision-maker. This separation allows us to study a focused question: \emph{which decisions can an LLM reliably make in an RPG, and how effectively does it self-correct when given feedback on invalid actions?}

We evaluate models at multiple scales, from 9B to 397B parameters, to identify the minimum model size required for consistent gameplay.

Our contributions are as follows:
\begin{itemize}
  \item We present a multi-agent system that addresses all facets of a retro RPG simultaneously --- exploration, combat, team building, resource management, and quest tracking --- through natural language reasoning, with each aspect handled by a specialized agent operating on structured game state. Our current evaluation covers the game up to the S.S.\ Anne, with full-game completion left as future work.
  \item We propose a multi-agent architecture with explicit separation between strategic LLM reasoning and mechanical game execution, using feedback-based validation that allows the LLM to self-correct.
  \item We evaluate models at multiple scales (9B to 397B parameters) and identify the minimum model size required for coherent gameplay in each agent role, demonstrating that complex RPG behavior is achievable through prompting alone, without reinforcement learning or fine-tuning.
  \item We provide an analysis of the decision boundaries between LLM reasoning and mechanical automation, contributing to the broader understanding of where LLM agents succeed and fail in complex interactive environments.
\end{itemize}

% --------------------------------------------------------------------------- %

\section{Related Work}
\label{sec:related-work}

Our work sits at the intersection of three active research areas: reinforcement learning for game playing, LLM-based game agents, and structured reasoning frameworks for language models.

\subsection{Reinforcement Learning for Game Playing}

Reinforcement learning has been the dominant paradigm for building game-playing AI. Deep Q-Networks demonstrated that a single architecture could learn to play dozens of Atari games from raw pixels \citep{mnih2015human}, while AlphaGo showed that RL combined with tree search could surpass human experts in Go \citep{silver2016mastering}. More recently, \citet{pleines2025pokemon} applied RL directly to Pok\'{e}mon Red, training agents over tens of thousands of hours of gameplay to navigate the game's opening areas. Their work demonstrates that RL can eventually produce viable game-playing behavior, but at enormous computational cost and with policies that remain brittle outside their training distribution. In contrast, our approach requires no training at all --- the LLM generalizes from its pre-training knowledge to make decisions in a zero-shot fashion.

\subsection{LLM-Based Game Agents}

A growing body of work explores LLMs as decision-making engines in interactive environments. Voyager \citep{wang2023voyager} deployed GPT-4 as a lifelong learning agent in Minecraft, using an automatic curriculum, a skill library of executable code, and iterative prompting with environment feedback. While Voyager demonstrates impressive open-ended exploration, Minecraft's sandbox nature differs fundamentally from the structured, narrative-driven progression of an RPG.

GITM \citep{zhu2023gitm} also targets Minecraft but takes a different approach from Voyager: it uses an LLM to decompose complex goals into sub-goals hierarchically, then executes each sub-goal with low-level action controllers. This separation between high-level reasoning and low-level execution parallels our own architecture, though GITM operates in a sandbox environment without narrative progression. Our system faces the additional challenge of pursuing a fixed storyline with specific quest objectives rather than open-ended goals.

Closest to our work is Pok\'{e}LLMon \citep{hu2024pokellmon}, which achieves human-parity performance in competitive Pok\'{e}mon battles on Showdown. Pok\'{e}LLMon introduces in-context reinforcement learning from battle feedback, knowledge-augmented generation to reduce hallucination, and consistent action generation to prevent panic switching. However, Pok\'{e}LLMon operates exclusively in the competitive battling domain --- a turn-based strategy game with complete information about available moves and types. It does not address overworld navigation, resource management, team building through wild encounters, or any of the other RPG mechanics that constitute the full game experience. Our agent attempts to play the complete RPG, coordinating across all these interleaved skills simultaneously.

CRADLE \citep{tan2024cradle} takes a different approach by targeting General Computer Control: a multimodal agent that interacts with software through screenshots and keyboard/mouse actions. CRADLE has been applied to complex games including Red Dead Redemption 2 and Stardew Valley, demonstrating broad applicability. Our approach differs in using direct memory reads instead of visual input, which eliminates perception errors and provides access to game state not visible on screen.

\subsection{LLM Agents for Pok\'{e}mon Red}

Concurrently with our work, two major AI labs have used Pok\'{e}mon Red as a testbed for their flagship models. Google reported that Gemini 2.5 Pro completed the entire game in approximately 400~hours across two runs, with full milestone tracking from the first rival battle through the Elite Four \citep{gemini25}. Anthropic demonstrated Claude 3.7 Sonnet defeating three Gym Leaders as part of their extended thinking announcement \citep{anthropic2025pokemon}, and subsequent iterations (Opus 4.5) obtained all eight badges but remained unable to complete Victory Road's boulder puzzles.

Both efforts use vision-based multimodal approaches with proprietary frontier models. Our work differs in three key respects: (1)~we use \emph{direct memory reads} instead of visual input, eliminating perception errors entirely; (2)~we employ \emph{open-weight models} (Qwen 3.5/3.6, 9B--397B parameters) rather than closed-source frontier systems; and (3)~our \emph{multi-agent architecture} decomposes the problem into seven specialized agents rather than relying on a single monolithic model. These design choices target a different research question: while the concurrent work demonstrates impressive game completion with frontier models, we study how a structured agent architecture can enable coherent gameplay with smaller, openly available models.

\subsection{Structured Reasoning for LLM Agents}

Our agent's output format draws on the ReAct framework \citep{yao2023react}, which interleaves reasoning traces with actions in a synergistic loop. Our agents follow a similar structure: each produces a \texttt{THINK:} line followed by an \texttt{ACTION:} directive. However, because we use models with native thinking/reasoning mode \citep{wei2022chain}, the bulk of the reasoning occurs internally within the model's chain-of-thought before generating the response. The explicit \texttt{THINK:} line functions primarily as an observable justification rather than the mechanism through which the model reasons.

LLM+P \citep{liu2023llmp} demonstrates that LLMs perform poorly at low-level spatial planning but excel at translating high-level intentions into formal problem specifications that classical planners can solve. This finding directly motivates our architecture's core design: the LLM decides \emph{where to go} while A* handles \emph{how to get there}, delegating spatial planning to an algorithm purpose-built for it.

Several works have explored reflection in LLM agents. Reflexion \citep{shinn2023reflexion} introduces verbal reinforcement learning, where an agent reflects on failed attempts to improve on subsequent trials of the same task. Agent-Pro \citep{zhang2024agentpro} extends this with policy-level reflection to evolve agent behavior over time. Park et al.\ \citep{park2023generative} introduce a memory architecture for generative agents where observations are stored in a memory stream, periodically synthesized into higher-level reflections, and retrieved to inform future decisions. Our battle reflection mechanism follows a similar observe--reflect--store pattern applied to combat: the agent observes damage outcomes, extracts type effectiveness facts, and persists them for retrieval in future battles (Section~\ref{sec:reflection}). This is related to Pok\'{e}LLMon's in-context reinforcement learning from battle feedback \citep{hu2024pokellmon}, but differs in that the learned knowledge is stored as persistent memory rather than maintained within a single context window, allowing it to accumulate across the entire playthrough.

\subsection{Multi-Agent LLM Architectures}

Recent frameworks such as MetaGPT \citep{hong2024metagpt} and AutoGen \citep{wu2023autogen} coordinate multiple LLM agents through inter-agent communication: agents with different roles converse, debate, and negotiate to solve complex tasks collaboratively. These architectures are well suited for open-ended problems where the task decomposition itself is ambiguous and benefits from discussion --- such as software engineering or multi-step reasoning.

Our system takes a different approach. Rather than having agents communicate with each other, we use a state-driven dispatcher that invokes exactly one specialized agent at a time based on the current game state. This design is motivated by two domain-specific constraints. First, Pok\'{e}mon Red's game states are mutually exclusive by design: the player is either in battle, navigating the overworld, in a shop, or naming a Pok\'{e}mon --- never in two contexts simultaneously. There is no ambiguity about which agent should act, making inter-agent negotiation unnecessary. Second, each LLM call incurs multi-second inference latency; a collaborative architecture that requires multiple rounds of agent-to-agent conversation would multiply this cost significantly. The result is a system where each agent receives only the information relevant to its domain, produces a single decision, and yields control back to the dispatcher.

% --------------------------------------------------------------------------- %

\section{System Overview}
\label{sec:system-overview}

Our system connects a large language model to a Game Boy emulator, enabling the LLM to observe game state and issue strategic decisions that are translated into button presses. The architecture is organized around three layers: a \emph{perception layer} that extracts structured state from emulator memory, a \emph{decision layer} composed of specialized LLM agents, and an \emph{execution layer} that translates high-level decisions into low-level input sequences. A feedback-based validation mechanism sits between the decision and execution layers, allowing the LLM to self-correct when it attempts invalid actions.

\subsection{Emulator Interface}

We use PyBoy \citep{pyboy}, an open-source Game Boy emulator written in Python, as our execution environment. Rather than interpreting the game visually through screen captures, we read game state directly from the emulator's working RAM (WRAM) at documented memory addresses. The Pok\'{e}mon Red fan community has produced comprehensive maps of the game's memory layout, providing access to player position, map identity, party composition, battle state, inventory contents, and active dialogue --- all as structured numerical data.

This approach offers two advantages over vision-based alternatives. First, it eliminates the need for an image recognition pipeline, removing an entire class of perception errors. Second, it provides access to information that is not always visible on screen, such as exact HP values, move power points, or the species of a wild Pok\'{e}mon before the battle animation completes.

The emulator runs synchronously with the agent: after each decision, the execution layer sends a sequence of button presses and advances the emulator by a fixed number of frames before the next state read. During LLM inference --- which takes several seconds per decision --- a wrapper keeps the emulator rendering on the main thread while the model runs on a background thread, preventing the display from freezing.

\subsection{Multi-Agent Architecture}

Rather than using a single monolithic prompt, we decompose gameplay into seven specialized agents, each responsible for a distinct aspect of the game:

\begin{itemize}
  \item \textbf{Overworld Agent}: Decides where to go next by examining two separate lists extracted from the current map: available warps (exits to other areas) and nearby points of interest (POIs) --- NPCs, items, and tall grass patches ---. It also receives the party's health status, badge progress, available money, a history of recent actions, and its own reasoning trace from the previous decision. This last element provides short-term intentionality: if the agent entered a Pok\'{e}mon Center intending to heal, it can recall that intention when the map changes and follow through. Outputs a target warp or POI for the Navigator to reach via pathfinding.
  \item \textbf{Battle Agent}: Selects combat actions (attack, use item, switch Pok\'{e}mon) based on the current battle state, including both combatants' types, HP, available moves, and status conditions.
  \item \textbf{Scout Agent}: Evaluates whether a newly encountered wild Pok\'{e}mon should be caught or defeated, considering the current party composition and whether the species (or its evolution line) is already represented.
  \item \textbf{Move Agent}: When a Pok\'{e}mon tries to learn a new move but already knows four, decides whether to replace an existing move or skip learning. Evaluates move stats, types, and the team's current coverage.
  \item \textbf{Shop Agent}: Decides what items to purchase at Pok\'{e} Marts based on current inventory, available funds, and party needs.
  \item \textbf{Lead Agent}: After each battle, evaluates the party and decides which Pok\'{e}mon should lead the team for the next encounter, considering HP status, level distribution, and training needs.
  \item \textbf{Naming Agent}: Generates creative nicknames for newly caught Pok\'{e}mon. While this agent does not affect gameplay outcomes, it illustrates the LLM's ability to perform creative tasks within the same architecture used for strategic reasoning.
\end{itemize}

Each agent follows the same interaction pattern: a set of state readers extract relevant information from WRAM, Jinja2 templates render this state into a natural language prompt, the LLM produces a response containing a \texttt{THINK:} reasoning trace followed by an \texttt{ACTION:} directive, and a parser extracts the structured decision. This output format is inspired by the ReAct framework \citep{yao2023react}, though the primary reasoning occurs within the model's native thinking mode rather than in the visible \texttt{THINK:} trace.

The main game loop dispatches agents based on the current game state, with a fixed priority order: battle handling takes precedence over dialogue, which takes precedence over overworld navigation. This ensures that interrupts --- such as a wild encounter during exploration --- are handled before the overworld agent attempts to continue navigating.

\subsection{Battle Reflection}
\label{sec:reflection}

After each battle, a reflection step analyzes the turn-by-turn history to extract type effectiveness facts --- for example, that Water moves are super effective against Rock types, or that Normal moves have no effect on Ghost types. These facts are persisted to a memory file and injected into future battle prompts, giving the battle agent an accumulating knowledge base built entirely from its own experience.

This mechanism is notable because the type effectiveness information is already encoded in the game's mechanics, but the LLM has no direct access to it. Instead of providing a complete type chart in the prompt, the agent learns these relationships incrementally through observation: when a move deals unexpectedly high or low damage, the battle reflect agent identifies the pattern and records it. Over time, this produces a personalized subset of the type chart grounded in encounters the agent has actually experienced.

\subsection{Decision vs. Execution Separation}

A central design principle is the strict separation between what the LLM decides and how those decisions are mechanically executed. The LLM outputs high-level intentions --- ``use Bubble'', ``go to the Pok\'{e} Center'', ``throw a Pok\'{e} Ball'' --- while dedicated driver modules handle the mechanical task of translating these into sequences of button presses:

\begin{itemize}
  \item A \textbf{Battle Menu Driver} navigates the in-game battle menu to select the correct move, open the bag to use an item, or switch to a specific party member.
  \item A \textbf{Navigator} computes A* paths between the player's current position and a target POI, then executes the walk step by step, handling collision and interrupting if an event (battle, dialogue) occurs mid-path.
  \item A \textbf{Name Entry} module translates a string into cursor movements on the in-game naming keyboard.
  \item A \textbf{Dialogue Advancer} detects and scrolls through text boxes, recording dialogue content for context while pressing the advance button at appropriate timing.
\end{itemize}

This separation is motivated by the observation that menu navigation and pathfinding require no intelligence --- they are deterministic translations of a goal into a sequence of inputs. Delegating these to code frees the LLM's limited context window and inference budget for decisions that genuinely benefit from reasoning.

\subsection{Action Validation and Feedback}

When the LLM attempts an action that is illegal or impossible given the current game state --- such as using a move with no remaining power points, selecting a fainted Pok\'{e}mon, or targeting a POI that does not exist on the current map --- the system provides explicit feedback explaining why the action failed and requests a new decision.

This design preserves the principle that the LLM is the sole decision-maker: the system constrains \emph{what information the agent receives}, not \emph{what the agent decides}. The feedback loop also produces a natural signal for evaluating agent competence --- each rejection represents a decision the LLM got wrong. We instrument this loop with a counter that tracks the frequency and type of invalid actions per agent, providing infrastructure for future quantitative analysis of decision quality.

In practice, we observe that the LLM typically self-corrects within one or two feedback rounds. The most common invalid actions include selecting POIs that do not exist in the current map and choosing moves that are not available to the active Pok\'{e}mon. If an agent exhausts its retry attempts without producing a valid action, the system cannot proceed autonomously --- this constitutes a termination condition for experimental runs (Section~\ref{sec:experiments}).

% --------------------------------------------------------------------------- %

\section{Agent Design}
\label{sec:agent-design}

This section details the prompt design and decision logic for each specialized agent. All agents share a common interaction pattern: they receive structured game state rendered through Jinja2 templates, produce a \texttt{THINK:} trace with brief reasoning, and output an \texttt{ACTION:} (or equivalent) directive that is parsed into a structured decision.

\subsection{Overworld Agent}

The overworld agent is the most frequently invoked and the most complex, responsible for all navigation decisions outside of battle. Its prompt is split into two parts: a system prompt with rules and the action format, and a state prompt with the current game context.

The state prompt presents two separate lists: \emph{interactables} (NPCs, items, tall grass patches) accessible via \texttt{INTERACT:<number>}, and \emph{warps} (exits to other maps) accessible via \texttt{GO:<number>}. Each entry includes an icon, a descriptive label, and the distance in steps. Interactables additionally carry a \emph{visit counter} showing how many times the agent has recently interacted with that POI. The full state includes:

\begin{itemize}
  \item The current map name.
  \item The party roster with nicknames, species, levels, and HP status (with visual indicators for low HP and fainted members).
  \item Pok\'{e} Ball count with a low-stock indicator, and current money.
  \item Badge progress (obtained vs remaining, with gym leader and city names).
  \item Key items currently in the inventory (e.g., Oak's Parcel), giving the agent awareness of quest-relevant objects it is carrying and can deliver.
  \item A rolling history of recent actions, providing short-term memory of previous decisions.
  \item The agent's own \texttt{THINK:} trace from the previous decision, providing short-term intentionality. Because each LLM call is stateless, the agent has no memory of what it was doing a moment ago. By feeding back its last reasoning trace, the agent can maintain coherent multi-step plans: if it entered a Pok\'{e}mon Center intending to heal, it sees that intention when the interior map loads and follows through rather than wandering off.
\end{itemize}

The system prompt provides minimal strategic guidance: it explains the two action types, notes that visited interactables are unlikely to yield new information, and informs the agent that tall grass contains wild Pok\'{e}mon whose battles grant experience. These are factual descriptions of game mechanics, not strategic directives.

The visit counter is a key design choice: by showing the agent how many times it has already interacted with each POI, we provide the information needed to avoid repetitive loops. The agent can reason for itself that visiting the same NPC five times is unproductive. Similarly, the rolling action history allows the agent to detect when it is cycling between the same locations and adjust its strategy accordingly.

\subsection{Battle Agent}

The battle agent handles all combat decisions. Its design is distinguished by a \emph{goal-based prompt switching} mechanism: the agent receives an entirely different system prompt depending on whether the battle goal is \textsc{defeat} or \textsc{catch}, as determined by the scout agent for wild encounters.

In \textsc{defeat} mode, the prompt provides game knowledge --- type advantage relationships and the STAB mechanic --- and notes that switching costs a turn, but leaves the strategic decision to the agent. In \textsc{catch} mode, the prompt explains the catch mechanics (lower HP increases catch rate, Pok\'{e} Balls are available as an action) and warns that STAB moves deal more damage than their base power suggests, but the agent decides when to weaken and when to throw.

The battle state provided to the agent includes:

\begin{itemize}
  \item The enemy Pok\'{e}mon's name, types, level, HP percentage, and any status conditions.
  \item The active Pok\'{e}mon's name, types, level, HP, and available moves --- each annotated with power, accuracy, PP remaining, type, physical/special category, a brief description, and a STAB indicator. Moves that are currently disabled (via the DISABLE status effect) are filtered from the available action space, preventing the agent from selecting them.
  \item Available party members for switching, with their types and levels.
  \item A turn-by-turn history tracking damage dealt and received.
  \item The available action space, which varies by context: \texttt{MOVE:<name>}, \texttt{SWITCH:<nickname>}, \texttt{BALL} (catch mode only), and \texttt{RUN} (wild battles only).
\end{itemize}

When a Pok\'{e}mon faints mid-battle, a dedicated \emph{switch-in decision} selects a replacement based on type matchups against the current enemy.
\subsection{Scout Agent}

Upon encountering a wild Pok\'{e}mon, the scout agent decides whether the battle goal should be \textsc{catch} or \textsc{defeat}. This decision runs before the battle agent is invoked, setting the strategic frame for the entire encounter.

The scout receives the wild Pok\'{e}mon's species, types, level, and full evolution line, alongside the current team roster with types and evolution lines. The system pre-computes several signals to assist the decision without dictating it:

\begin{itemize}
  \item A summary of types already covered by the team.
  \item A warning if the wild Pok\'{e}mon's evolution line is already represented on the team.
  \item A warning if the wild Pok\'{e}mon's types overlap with existing team members.
  \item A highlight if the wild Pok\'{e}mon would add types not yet covered.
\end{itemize}

These pre-computed signals reduce the reasoning burden on the LLM --- it does not need to manually cross-reference types across the team --- while leaving the final catch-or-defeat decision entirely to the agent. The prompt establishes a priority hierarchy (type diversity over evolution potential over avoiding duplicates) but the agent weighs these factors in context.

\subsection{Move Agent}

When a Pok\'{e}mon attempts to learn a new move but already knows four, the move agent decides whether to replace an existing move or skip learning. The agent receives the Pok\'{e}mon's species, the new move's stats (power, accuracy, PP, type, category, description), and the same details for all four current moves.

The prompt provides general principles --- higher power is generally better, type coverage across different types is valuable, keeping two moves of the same type is usually redundant, and status moves are situational. The agent outputs \texttt{FORGET <move name>} to replace a specific move or \texttt{KEEP} to skip learning.

\subsection{Shop Agent}

The shop agent is invoked when the player enters a Pok\'{e} Mart. It receives the list of available items with prices, the player's current money, Pok\'{e} Ball count, and party state. The agent outputs a purchase list in the format \texttt{BUY: <item\_number> x<quantity>}, or \texttt{NONE} if no purchases are needed.

The prompt provides minimal guidance: it instructs the agent to buy what it considers necessary based on team and inventory needs, to keep some money as reserve, and to buy nothing if nothing is needed. Beyond these general principles, the specific purchasing decisions are left to the agent.

\subsection{Lead Agent}

After each battle, the lead agent evaluates the party and decides which Pok\'{e}mon should be in the first position for the next encounter. The agent receives the full party roster with each member's nickname, species, level, current HP, and health status.

The prompt offers soft considerations rather than strict rules: a Pok\'{e}mon with very low HP should probably not lead, members that could benefit from experience are worth considering for the front position, and if the current lead is healthy then keeping it is often reasonable. The agent weighs these factors and decides for itself.

The agent outputs \texttt{KEEP} to retain the current lead or \texttt{SWITCH:<number>} to promote a different party member.

\subsection{Naming Agent}

The naming agent generates nicknames for newly caught Pok\'{e}mon. It receives the species name, the full evolution line, and the current team's nicknames to avoid duplicates. The prompt instructs the agent to choose a name that remains fitting across all evolutionary stages --- for instance, a name for a Charmander should still suit a Charizard.

The output is constrained to 10 uppercase characters, matching Pok\'{e}mon Red's character limit. A post-processing step strips any characters outside the game's allowed character set and truncates names that exceed the limit. As noted in Section~\ref{sec:system-overview}, this agent does not influence gameplay outcomes --- it exists purely to make each playthrough more entertaining to observe.

% --------------------------------------------------------------------------- %

\section{Mechanical Automation Layer}
\label{sec:automation}

While the LLM handles all strategic decisions, a set of automated modules handle the mechanical translation of those decisions into emulator inputs. These components perform tasks that are fully deterministic and require no reasoning --- they are the equivalent of a human player's muscle memory for pressing buttons.

\subsection{Pathfinding and Navigation}

When the overworld agent selects a POI, the navigator module computes a path using the A* algorithm over a 2D walkability grid of the current map. A* uses Manhattan distance as its heuristic and accounts for dynamically blocked tiles such as NPCs. The navigator then executes the path step by step, pressing directional buttons and advancing frames between each tile.

Navigation is designed for interruption safety. At each step, the navigator checks whether a battle has started, dialogue has appeared, or the player has warped to a new map. If a battle interrupts navigation, the pending POI is preserved and navigation resumes automatically after the battle is resolved, since the original destination was never reached. This is important for maintaining coherence: without resumption, the agent would observe that its last decision produced no visible result --- it chose to go to a POI but arrived somewhere unexpected --- leading to confusion in subsequent reasoning. By completing the interrupted action, the system ensures that the agent's decisions have the outcomes it intended. A map transition triggers a wait period for the game to stabilize before continuing.

The path is recalculated before each step rather than computed once upfront, because NPCs move between tiles and can block previously valid routes. The walkability grid itself includes manual overrides for cases where the ROM's collision data does not match actual walkability --- for example, certain cave wall tiles in Mt.\ Moon that the game marks as walkable but are in fact blocked. These overrides are defined as a static table of tile corrections.

The pathfinder also handles directional ledges --- tiles that can only be traversed in one direction (e.g., jumping down a ledge but not climbing back up). These are modeled as one-way edges in the walkability grid, allowing A* to find valid paths that respect the game's movement constraints.

Different POI types require different approach strategies: warps activate by walking onto the tile, NPCs require facing the target and pressing A, and map connections require walking to the edge and stepping off.

\subsection{Battle Menu Driver}

The battle menu driver translates the LLM's decisions into the specific button sequences needed to navigate Pok\'{e}mon Red's nested menu system. When the battle agent outputs \texttt{MOVE:Bubble}, the driver reads the current cursor position from memory, calculates the minimal sequence of directional presses to reach the correct move slot, and confirms the selection.

The same mechanism handles item usage, Pok\'{e} Ball throwing, and party switching --- each requiring navigation through different sub-menus. By tracking the cursor position in WRAM, the driver minimizes unnecessary button presses.

\subsection{Dialogue Advancement}

A dialogue advancer detects when a text box appears on screen by reading the tilemap from WRAM. It waits for the letter-by-letter text animation to complete, captures the text content, and presses A to advance. When text scrolls, the module detects overlap with previously captured lines and extracts only the new content.

The captured dialogue serves a dual purpose: it drives the game forward mechanically, and it provides context to the overworld agent's action history, helping it reason about what has happened recently. An NPC saying ``Please deliver this parcel'' is both a text box to dismiss and contextual information that informs future navigation decisions.

The advancer stops based on configurable conditions: the appearance of a battle menu, a YES/NO choice prompt, a naming keyboard, or a prolonged silence indicating a scene transition.

\subsection{Event Processing}

An event processor sits between the navigator and the game, handling contextual interactions that arise during gameplay. When the player interacts with a Pok\'{e} Center nurse and a YES/NO prompt appears, the event processor automatically accepts. When a shop menu appears, it delegates to the shop agent. When a naming keyboard appears after catching a Pok\'{e}mon, it delegates to the naming agent.

This module acts as a dispatcher: it detects the type of event from WRAM state and routes it to the appropriate handler, whether that is an automated response (accepting healing) or an LLM agent (deciding what to buy). When the agent encounters an unrecognized menu, the event processor handles it based on context: PC-related menus are dismissed, while other unknown menus default to selecting the first option. This prevents the agent from getting stuck in interfaces it was not designed to handle.

\subsection{Name Entry}

When the naming agent produces a nickname, a name entry module translates the string into cursor movements on the in-game keyboard. The keyboard is a 9-column, 3-row grid of uppercase characters. For each character, the module calculates the row and column offset from the current cursor position, issues the corresponding directional presses, and confirms with A. Entry is finalized by pressing START.

\subsection{Design Rationale}

Every component in this layer shares a common property: given the same input, it will always produce the same output. There is no decision to be made about \emph{how} to walk from tile A to tile B, or \emph{how} to press down twice to reach the third move in a menu. Automating these tasks is not an assistance to the LLM --- it is the removal of noise that would otherwise consume tokens on trivially deterministic problems, freeing the model's limited context window for decisions that genuinely benefit from reasoning. This division follows the principle established by LLM+P \citep{liu2023llmp}: LLMs are effective at formulating intentions but unreliable at low-level spatial planning, which is better delegated to classical algorithms.

% --------------------------------------------------------------------------- %

\section{Experiments}
\label{sec:experiments}

\subsection{Setup}

All experiments are conducted on a MacBook Pro 14-inch (November 2024) with an Apple M4 Pro chip and 24~GB of unified memory. The game runs on PyBoy, an open-source Game Boy emulator written in Python, at normal emulation speed. Each run starts from the beginning of the game (New Game) with no prior state. The only human inputs are the player and rival names, entered manually before the agent takes control; all decisions from the player's bedroom onward are made entirely by the LLM. All models are served via cloud inference through OpenRouter, with thinking/reasoning mode enabled.

We evaluate six models from the Qwen family \citep{qwen3}, spanning dense and mixture-of-experts architectures:

\begin{itemize}
  \item \textbf{Qwen 3.5 9B} --- dense, 9B parameters.
  \item \textbf{Qwen 3.5 27B} --- dense, 27B parameters.
  \item \textbf{Qwen 3.5 35B-A3B} --- mixture-of-experts, 35B total / 3B active parameters.
  \item \textbf{Qwen 3.5 122B-A10B} --- mixture-of-experts, 122B total / 10B active parameters.
  \item \textbf{Qwen 3.5 397B-A17B} --- mixture-of-experts, 397B total / 17B active parameters.
  \item \textbf{Qwen 3.6 Plus} --- latest generation, served as a managed endpoint.
\end{itemize}

For each model, all seven agents use the same model. We conduct a single run per model. Each run terminates under one of three conditions:

\begin{enumerate}
  \item \textbf{Completion}: The agent reaches the S.S.\ Anne.
  \item \textbf{Retry exhaustion}: Any agent fails to produce a valid action after its maximum retry attempts, requiring human intervention and breaking autonomy.
  \item \textbf{Unrecoverable loop}: The agent enters a repetitive cycle that fills its entire action history without making progress.
\end{enumerate}

\subsection{Metrics}

\subsubsection{Per-Agent Decision Metrics}

We instrument each agent to continuously log decision statistics:

\begin{itemize}
  \item \textbf{Overworld}: Total decisions, broken down by type (warp, interact, tall grass).
  \item \textbf{Battle}: Total decisions by type (attack, switch, ball, run). Wins, losses (whiteouts), and captures.
  \item \textbf{Battle Reflect}: Total lessons recorded. Accuracy of learned type effectiveness compared against the game's actual type chart.
  \item \textbf{Scout}: Decisions to catch vs.\ defeat.
  \item \textbf{Lead}: Decisions to swap the lead vs.\ keep the current one.
  \item \textbf{Shop}: Number of purchase transactions.
  \item \textbf{Move Learning}: Decisions to learn a new move vs.\ keep the existing moveset.
  \item \textbf{Naming}: Number of nicknames generated.
  \item \textbf{All agents}: Total unparseable responses (responses that fail to match the expected output format).
\end{itemize}

\subsubsection{Narrative Progression Checkpoints}

We record the wall-clock time at which the agent reaches each story milestone, along with the team's average level and party size at that point:

\begin{enumerate}
  \item Obtain starter Pok\'{e}mon
  \item Deliver Oak's Parcel
  \item Reach Viridian Forest
  \item Defeat Brock (Badge 1)
  \item Reach Mt.\ Moon
  \item Reach Cerulean City
  \item Defeat Misty (Badge 2)
  \item Help Bill
  \item Reach Vermilion City
  \item Complete the S.S.\ Anne
\end{enumerate}

% --------------------------------------------------------------------------- %

\section{Results}
\label{sec:results}

\subsection{Summary}

Table~\ref{tab:summary} presents a high-level comparison of all evaluated models. For each run, we report the furthest narrative checkpoint reached, total wall-clock time, and key decision statistics across agents.

\begin{table}[h]
\centering
\small
\caption{Summary of experimental runs by model.}
\label{tab:summary}
\begin{tabular}{lrrrrrrrr}
\toprule
Model & Time & OW & Battle & Wins & Losses & Catches & Unparse. & Outcome \\
\midrule
Qwen 3.5 9B       & 0.1h &   6 &   --- & --- & --- & --- &  3 & Failed (unparseable) \\
Qwen 3.5 27B      &  30h & 790 & 1167  & 413 &   2 &  16 & 32 & Badge 2 (Misty) \\
Qwen 3.5 35B-A3B  & \multicolumn{8}{c}{\emph{--- pending ---}} \\
Qwen 3.5 122B-A10B & \multicolumn{8}{c}{\emph{--- pending ---}} \\
Qwen 3.5 397B-A17B & \multicolumn{8}{c}{\emph{--- pending ---}} \\
Qwen 3.6 Plus     &  37h & 894 & 1553  & 351 &  15 &   4 &  9 & Badge 2 + Bill \\
\bottomrule
\end{tabular}
\end{table}

\subsection{Narrative Progression}

Tables~\ref{tab:cp-27b} and~\ref{tab:cp-36plus} show the wall-clock time and party composition at each narrative checkpoint for the two models that achieved meaningful progress. Qwen 3.5 9B failed before reaching any checkpoint (Section~\ref{sec:obs}).

\begin{table}[h]
\centering
\small
\caption{Qwen 3.5 27B --- checkpoint progression.}
\label{tab:cp-27b}
\begin{tabular}{rlrl}
\toprule
\# & Checkpoint & Time & Party \\
\midrule
1 & Obtain starter    & 0h 14m  & Lv5 Bulbasaur \\
2 & Deliver Parcel    & 1h 03m  & Lv7 Bulbasaur \\
3 & Viridian Forest   & 2h 26m  & Lv9 Bulbasaur, Lv4 Rattata, Lv3 Spearow \\
4 & Defeat Brock      & 5h 02m  & Lv13 Bulbasaur, Lv9 Rattata, Lv5 Spearow, Lv3 Caterpie \\
5 & Reach Mt.\ Moon   & 6h 51m  & Lv14 Bulbasaur, Lv11 Spearow, Lv10 Rattata, Lv6 Caterpie \\
6 & Reach Cerulean    & 22h 26m & Lv24 Ivysaur, Lv16 Spearow, Lv16 Rattata, Lv11 Geodude, Lv11 Butterfree, Lv4 Pikachu \\
7 & Defeat Misty      & 22h 46m & Lv25 Ivysaur, Lv17 Spearow, Lv17 Rattata, Lv11 Geodude, Lv11 Butterfree, Lv4 Pikachu \\
\bottomrule
\end{tabular}
\end{table}

\begin{table}[h]
\centering
\small
\caption{Qwen 3.6 Plus --- checkpoint progression.}
\label{tab:cp-36plus}
\begin{tabular}{rlrl}
\toprule
\# & Checkpoint & Time & Party \\
\midrule
1 & Obtain starter    & 0h 15m  & Lv5 Bulbasaur \\
2 & Deliver Parcel    & 0h 28m  & Lv5 Bulbasaur \\
3 & Viridian Forest   & 0h 50m  & Lv6 Bulbasaur, Lv4 Pidgey \\
4 & Defeat Brock      & 17h 57m & Lv16 Pikachu, Lv14 Bulbasaur, Lv13 Pidgey, Lv7 Kakuna \\
5 & Reach Mt.\ Moon   & 18h 56m & Lv17 Pikachu, Lv15 Pidgey, Lv14 Bulbasaur, Lv8 Kakuna \\
6 & Reach Cerulean    & 28h 34m & Lv20 Pikachu, Lv18 Pidgeotto, Lv17 Ivysaur, Lv16 Geodude, Lv15 Beedrill \\
7 & Defeat Misty      & 28h 53m & Lv20 Pikachu, Lv18 Pidgeotto, Lv17 Ivysaur, Lv16 Geodude, Lv16 Beedrill \\
8 & Help Bill         & 34h 22m & Lv23 Pikachu, Lv22 Pidgeotto, Lv19 Geodude, Lv19 Ivysaur, Lv17 Beedrill \\
\bottomrule
\end{tabular}
\end{table}

% TODO: Add checkpoint tables for Qwen 3.5 35B-A3B, 122B-A10B, 397B-A17B

\subsection{Per-Agent Statistics}

Table~\ref{tab:agents} breaks down LLM invocations and key statistics per agent for each completed run.

\begin{table}[h]
\centering
\small
\caption{Per-agent decision statistics for completed runs.}
\label{tab:agents}
\begin{tabular}{llrrrr}
\toprule
Agent & Metric & 9B & 27B & 3.6 Plus \\
\midrule
Overworld      & Decisions      &   6 & 790 &  894 \\
               & Avg time (s)   & 32.7 & 8.3 & 19.2 \\
Battle         & Decisions      & --- & 1167 & 1553 \\
               & Avg time (s)   & --- & 24.6 & 27.3 \\
               & Attacks        & --- &  976 & 1130 \\
               & Switches       & --- &  124 &  318 \\
Battle Reflect & Facts learned  & --- &   23 &   30 \\
               & Incorrect      & --- &    3 &    5 \\
               & Avg time (s)   & --- & 32.3 & 32.0 \\
Scout          & Evaluated      & --- &  370 &  305 \\
               & Catches        & --- &   16 &    4 \\
               & Avg time (s)   & --- & 22.9 & 17.6 \\
Lead           & Decisions      & --- &  410 &  365 \\
               & Swaps          & --- &   30 &   34 \\
               & Avg time (s)   & --- & 20.7 & 21.4 \\
Shop           & Transactions   & --- &    3 &   10 \\
               & Avg time (s)   & --- & 30.6 & 36.6 \\
Naming         & Names          & --- &    6 &    5 \\
               & Avg time (s)   & --- & 331.1 & 105.7 \\
Move           & Decisions      & --- &    2 &  --- \\
               & Avg time (s)   & --- & 39.9 &  --- \\
\bottomrule
\end{tabular}
\end{table}

\subsection{Observations}
\label{sec:obs}

\textbf{Qwen 3.5 9B.} The smallest model failed immediately --- it produced 3 consecutive unparseable responses after only 6 overworld decisions and was unable to progress past the starting town. Thinking/reasoning mode did not compensate for the lack of capacity.

\textbf{Qwen 3.5 27B.} Reached Misty (Badge 2) in 30 hours. The agent chose Bulbasaur as its starter, but this did not translate into an immediate type advantage against Brock: Bulbasaur does not learn Vine Whip (its first Grass-type move) until level 14, and the agent challenged the gym before reaching that level. The victory came instead through correct use of Leech Seed, a strategy the agent discovered on its own. Mt.\ Moon remained the major bottleneck, taking over 15 hours (6.9h to 22.4h). The agent never entered Bill's house --- it rarely explored building interiors throughout the run, suggesting that entering buildings is a navigation pattern the model struggles to prioritize.

\textbf{Qwen 3.6 Plus.} Reached the furthest of all models tested so far, defeating Misty and completing Bill's quest (S.S.\ Ticket obtained) in 37 hours. The Brock bottleneck was severe: despite choosing Bulbasaur, the agent repeatedly challenged the gym but never managed to execute the Leech Seed strategy that the 27B model discovered. Instead, it spent nearly 18 hours grinding and retrying until Bulbasaur learned Vine Whip at level 14, brute-forcing the badge through raw stats rather than strategic play. Mt.\ Moon took approximately 10 hours (18.9h to 28.6h). The run terminated after 2.5 hours without progress past Cerulean City --- the exit to Route 5 requires passing through a building interior, a non-obvious navigation pattern the model could not infer.

\textbf{Mt.\ Moon as a universal bottleneck.} Across all models that reached Mt.\ Moon, the cave represented the single largest time sink. This is driven by two compounding factors. First, the cave has an extremely high wild encounter rate, and the agents rarely choose to run from battles unless the situation is clearly unfavorable. The cumulative effect of dozens of unplanned battles drains the party's HP and move PP, eventually forcing the agent to leave the cave to heal at a Pok\'{e}mon Center --- losing all navigation progress in the process. Second, upon re-entering, the agent has no memory of how far it previously progressed, and the cave's internal sections share similar names, causing confusion. In practice, all models required significant leveling before the party could survive the full cave traversal in a single attempt.

% TODO: Add observations for Qwen 3.5 35B-A3B, 122B-A10B, 397B-A17B

% --------------------------------------------------------------------------- %

\section{Discussion}
\label{sec:discussion}

\subsection{Where the LLM Succeeds}

The results reveal several decision domains where the LLM performs reliably well, even at moderate model sizes.

\textbf{Tactical battle decisions.} The battle agent consistently selects appropriate moves given type matchup information and move stats. The Qwen 3.5 27B run provides a striking example: the agent independently discovered that Leech Seed could defeat Brock's Rock-type team despite having no super-effective damage moves available, demonstrating genuine tactical reasoning rather than pattern matching. Across all runs, the battle agent maintained low unparseable response rates relative to the number of decisions made.

\textbf{Team composition.} The scout agent maintains type diversity effectively, avoiding duplicate evolution lines and favoring Pok\'{e}mon that add new type coverage. The pre-computed signals in the prompt (type overlap warnings, new type highlights) reduce the reasoning burden, but the catch-or-defeat decision itself remains with the agent and is generally sound.

\textbf{Learning from experience.} The battle reflection mechanism successfully extracts type effectiveness facts from combat, building a persistent knowledge base that improves future decisions. Across the 27B and 3.6 Plus runs, the agent learned 23 and 30 facts respectively, with relatively low error rates (3 and 5 incorrect facts). This demonstrates that an LLM can incrementally learn game mechanics through observation and apply that knowledge in subsequent encounters.

\textbf{Decisions over explicit option sets.} A common pattern emerges: the LLM excels when presented with a closed list of options and sufficient context to evaluate them. Choosing a move from four options, selecting a party lead from six members, or deciding between catch and defeat --- these are decisions where the action space is small, fully enumerated, and accompanied by relevant information. The structured prompt format plays a key role here: by presenting options as numbered lists with associated data, we reduce the task to selection rather than generation.

\subsection{Where the LLM Fails}

Several failure modes recur across models and runs, pointing to fundamental limitations of stateless LLM decision-making in extended gameplay.

\textbf{Spatial navigation without persistent memory.} The Mt.\ Moon bottleneck illustrates the core problem: the agent cannot maintain a navigation plan across multiple LLM calls. Each decision is made independently, with only a rolling history of recent actions and the previous reasoning trace as context. When the agent exits the cave to heal and re-enters, it has no record of how far it previously progressed. The result is repeated partial traversals of the same sections, with progress gated by the party's ability to survive the entire cave in a single attempt rather than by navigation skill.

\textbf{Inferring non-obvious routes.} The Cerulean City blocker reveals a harder problem: the exit to Route 5 requires entering a building and walking through to the other side. This navigation pattern is never presented as an explicit option --- the agent sees a building entrance as a warp like any other, with no indication that it leads to a different route. Inferring this requires world knowledge about Pok\'{e}mon Red's map layout that is not available in the prompt, and the agent cannot discover it through exploration alone because it has no reason to enter an arbitrary building and walk to its far exit.

\textbf{Under-exploration of building interiors.} Related to the above, the 27B agent rarely entered buildings throughout its run --- most notably, it never visited Bill's house despite reaching Cerulean City. Warp options to building interiors compete with warps to new routes, and the agent systematically prefers the latter. This suggests that the overworld prompt's framing, which encourages forward progression, may inadvertently discourage the exploratory behavior needed to discover quest-relevant interiors.

\textbf{Short-term navigation loops.} Agents occasionally cycle between the same warps, which unlike interactable POIs do not carry a visit counter. However, these loops are typically short-lived: the rolling action history provides enough context for the model to recognize the repetition and break out after a few iterations.

\subsection{Effect of Model Scale}

The three completed runs suggest a non-linear relationship between model size and gameplay quality.

At the smallest scale, Qwen 3.5 9B could not reliably produce parseable output, failing after only 6 overworld decisions. This establishes a clear minimum threshold: models below approximately 10B parameters lack the capacity to follow the structured output format consistently, regardless of the quality of their reasoning.

Qwen 3.5 27B and Qwen 3.6 Plus both reached Badge 2, but through qualitatively different play styles. The 27B model defeated Brock in 5 hours through a creative strategy (Leech Seed), while the larger 3.6 Plus model spent 18 hours brute-forcing the same gym through grinding. This counterintuitive result --- a smaller model solving a problem more elegantly --- suggests that model scale does not uniformly improve strategic reasoning. The 3.6 Plus model did ultimately progress further (completing Bill's quest), indicating that scale may matter more for sustained coherence over long runs than for individual tactical decisions.

The pending runs on Qwen 3.5 35B-A3B, 122B-A10B, and 397B-A17B will provide additional data points to characterize this relationship more precisely, particularly regarding whether mixture-of-experts architectures with fewer active parameters behave more like their total or active parameter count.

\subsection{Boundaries of the Decision/Execution Separation}

Our architecture draws a strict line between LLM reasoning and mechanical execution. The results validate this separation in most cases: pathfinding, menu navigation, dialogue advancement, and name entry all function reliably as deterministic translations of high-level intentions into button presses. The LLM never needs to reason about \emph{how} to walk to a tile or \emph{how} to select a menu option, and this frees its context window for genuinely strategic decisions.

However, the separation has limits. The execution layer filters the decision space presented to the LLM: only POIs reachable by A* appear as options, and only moves with remaining PP are shown. This filtering is generally beneficial --- it prevents the agent from attempting impossible actions --- but it also means the agent never sees what it \emph{cannot} do. In many cases this mirrors the human experience: a player in Cerulean City also cannot see the exit to Route 5 until they discover the building that connects to it. But in other cases the gap is meaningful: a human player can visually spot an item behind a cuttable tree and reason that they need HM Cut to reach it, while our agent has no awareness that the item exists at all. The Cerulean City blocker is a direct consequence of this design: the exit to Route 5 is not reachable by A* unless the agent first passes through a building that unblocks the path, so the warp never appears in the option list. The agent cannot reason about a destination it has never been shown.

\subsection{Comparison with Concurrent Work}

Concurrently with our work, two major AI labs have demonstrated frontier models playing Pok\'{e}mon Red. Gemini 2.5 Pro completed the entire game in approximately 400~hours \citep{gemini25}, and Claude progressed from three gym badges (3.7 Sonnet) to all eight badges (Opus 4.5), though it was unable to complete Victory Road \citep{anthropic2025pokemon}. Both efforts use vision-based multimodal approaches with proprietary models whose parameter counts and architectures are undisclosed.

These results establish that sufficiently powerful frontier models can complete the game when given enough time and multimodal perception. Our work asks a complementary question: \emph{can a structured multi-agent architecture achieve coherent RPG gameplay with open-weight models at a fraction of the scale?} The answer, based on our experiments, is a qualified yes --- a 27B-parameter model reaches Badge~2 in under 30~hours.

Three properties distinguish our approach from the frontier model demonstrations. First, \textbf{openness}: our system uses open-weight models, documented WRAM addresses, and an open-source emulator --- every component can be inspected, modified, and fine-tuned by third parties. The Gemini and Claude runs use proprietary models accessible only through APIs, limiting the ability to adapt or build upon their approach. Second, \textbf{interpretability}: because each decision is made by a specialized agent with a domain-specific prompt, the reasoning behind any individual choice can be traced to its agent, its input state, and its output format. The vision-based approaches prioritize generality --- requiring no game-specific memory reads --- but offer less granular insight into individual decisions. Third, \textbf{efficiency}: while direct comparison is imperfect given the differences in game progress, our 27B model reaches Badge~2 in 23~hours --- suggesting that structured architecture can compensate for raw model scale when the decision space is well-decomposed.

% --------------------------------------------------------------------------- %

\section{Limitations and Future Work}
\label{sec:limitations}

\subsection{Manual Interventions}

During the experimental runs, two edge cases in the mechanical execution layer required brief human intervention. Neither situation affected gameplay progress or the LLM's decision quality --- both are straightforward to patch in the execution layer and reflect gaps in edge case coverage rather than limitations of the agent architecture.

\begin{enumerate}
  \item The multi-turn trapping moves Wrap (Ekans) and Bind (Onix) lock both combatants for several turns during which no action can be taken. The battle menu driver was not designed for this state and in one instance navigated into the item menu, requiring manual recovery.
  \item After helping Bill, interacting with his PC opens a Pok\'{e}dex entry screen that was not handled by the event processor, requiring the menu to be dismissed manually.
\end{enumerate}

\subsection{Building Interior Exploration}

The most immediate barrier to further progress is the agent's tendency to avoid building interiors. The exit from Cerulean City to Route 5 --- the next step in the game's narrative --- requires entering a building and walking through to the other side, a pattern that no model discovered autonomously. More broadly, quest-relevant locations such as Bill's house are inside buildings that the agent rarely enters.

Two complementary approaches could address this. First, prompt-level changes could encourage interior exploration more explicitly --- for instance, by highlighting warps that have never been taken, giving the agent a signal that unexplored interiors may contain something valuable. Second, a richer warp annotation system could provide the agent with information about what lies behind each warp: whether the destination has been visited before, and if so, what POIs and exits it contains. This would allow the agent to reason about whether entering a building is worthwhile without having to explore it blindly. Such annotations could also improve Mt.\ Moon navigation, where the agent currently has no way to distinguish between cave exits it has already explored and those that lead to new areas.

\subsection{Fixed Party of Six}

The current system has no mechanism for managing the PC storage system. Once the party reaches six members, any additional captures are sent to the PC box and become inaccessible. This means the agent's team composition is effectively locked after the sixth catch --- it cannot swap in a stronger Pok\'{e}mon caught later or retrieve a specialist for a difficult gym. A natural extension would be a dedicated \emph{PC Agent} that manages team composition by depositing and withdrawing Pok\'{e}mon from storage, activated when the overworld agent interacts with the PC in a Pok\'{e}mon Center, following the same delegation pattern used for the shop agent when the player interacts with a mart clerk.

\subsection{Single Run per Model}

Each model is evaluated with a single run, which limits our ability to measure variance. The agent's trajectory is highly sensitive to early decisions --- choosing Charmander instead of Bulbasaur as a starter, for instance, drastically changes the Brock gym difficulty. A single run cannot distinguish between a model's general capability and a lucky or unlucky sequence of choices. Future work should include multiple independent runs per model to establish confidence intervals on progression speed and narrative reach.

\subsection{No Overworld HM Usage}

Pok\'{e}mon Red requires Hidden Machine (HM) moves --- Cut, Surf, Strength, and Fly --- to access areas blocked by trees, water, or boulders, and to fast-travel between cities. The current system has no mechanism for using moves outside of battle, which places a hard ceiling on how far the agent can progress. (Flash, the fifth HM, is not a blocker: it only reduces screen visibility in certain caves, but since our agent reads game state from WRAM rather than from the screen, darkness has no effect on its perception or navigation.)

Implementing HM support involves two shared challenges and four HM-specific designs. The shared challenges are: (1)~deciding which party member should learn the HM move, which could be triggered by the overworld agent recognizing the relevant key item and delegating to the existing move learning agent; and (2)~invoking the move through the overworld menu rather than in battle. The HM-specific designs vary in complexity:

\begin{itemize}
  \item \textbf{Cut}: Cuttable trees can be marked as interactable POIs; when the agent selects one, the execution layer opens the party menu and uses Cut automatically.
  \item \textbf{Fly}: The overworld agent receives a list of previously visited cities as fly destinations, choosing one like any other navigation target.
  \item \textbf{Surf}: Water tiles are reclassified as walkable when a party member knows Surf, making coastal POIs and warps reachable by A*. The execution layer detects the water boundary and triggers Surf at the appropriate moment.
  \item \textbf{Strength}: This is the most challenging HM to support, as it requires solving boulder puzzles --- pushing rocks in specific sequences to clear a path. Unlike the other HMs, this demands spatial reasoning over multiple steps and would require a dedicated approach to represent the puzzle state and guide the agent through the solution.
\end{itemize}

\subsection{Dependence on Documented Memory Layout}

Our entire perception layer relies on WRAM addresses mapped by the Pok\'{e}mon Red fan community. This provides precise, noise-free access to game state but does not generalize to other games without equivalent documentation. Extending the system to other retro RPGs would require a comparable reverse-engineering effort for each title's memory layout.

\subsection{Fine-Tuning Smaller Models}

Our evaluation uses general-purpose models with no task-specific training. The structured logs produced by each run --- thousands of (game state, reasoning trace, action) tuples per agent --- constitute a natural fine-tuning dataset. Training a smaller model on successful decisions from a larger model's run could yield an agent that matches the larger model's gameplay quality at a fraction of the inference cost and latency. This distillation approach is particularly promising for agents with highly repetitive decision patterns, such as battle and lead selection, where the action space is small and the correct decisions are well-represented in the logs.

\subsection{Single Model for All Agents}

In our experiments, all seven agents use the same model. However, agents differ dramatically in complexity: the naming agent generates a creative string from a short prompt, while the overworld agent reasons over maps, history, and party state. Using a smaller, faster model for simpler agents (naming, shop, move learning) and reserving a larger model for the overworld and battle agents could significantly reduce inference cost and latency without sacrificing decision quality. This heterogeneous configuration is straightforward to implement --- the architecture already treats each agent as an independent consumer of an LLM interface --- but remains unexplored in our current evaluation.

% --------------------------------------------------------------------------- %

\section{Conclusion}
\label{sec:conclusion}

We presented an autonomous agent that plays Pok\'{e}mon Red using a large language model as its sole decision-making engine, with no reinforcement learning, fine-tuning, or task-specific training. The system decomposes gameplay into seven specialized agents --- overworld navigation, battle strategy, catch-or-defeat evaluation, move learning, shop optimization, party lead selection, and nickname generation --- each operating on structured prompts constructed from emulator memory reads, with a battle reflection mechanism that builds persistent type effectiveness knowledge from combat experience.

Our experiments across models ranging from 9B to 397B parameters reveal that a 27B-parameter model is sufficient to play the game's opening act coherently, defeating two gym leaders and managing a diverse team of six Pok\'{e}mon. The results highlight a clear separation between decision types: the LLM excels at choices presented as selection over explicit, well-annotated options --- picking a move in battle, evaluating whether to catch a wild Pok\'{e}mon, choosing a party lead --- but struggles with spatial navigation that requires persistent memory across stateless calls and with discovering non-obvious routes through building interiors.

A notable finding is that larger models do not uniformly produce better gameplay. Qwen 3.5 27B defeated Brock through a creative Leech Seed strategy, while the larger Qwen 3.6 Plus spent 18 hours brute-forcing the same gym through grinding. Scale appears to matter more for sustained coherence over long runs than for individual tactical decisions.

The architecture's strict separation between strategic reasoning and mechanical execution proved effective: deterministic modules for pathfinding, menu navigation, and dialogue advancement eliminated an entire class of errors while freeing the LLM's context window for decisions that genuinely benefit from reasoning. The remaining barriers to further progression --- building interior exploration, HM usage, and PC-based team management --- are all addressable through extensions to the existing architecture rather than fundamental redesigns.

Recent demonstrations by Google and Anthropic --- completing or nearly completing Pok\'{e}mon Red with Gemini 2.5 Pro and Claude Opus 4.5 respectively --- validate the game as a meaningful benchmark for LLM agents. Our work complements these efforts by showing that structured agent architecture can enable coherent gameplay with open-weight models at a fraction of the scale. Where these demonstrations showcase the capabilities of frontier models, our system provides an interpretable framework for studying how each gameplay decision is made, built entirely on open-weight models and open-source tools.

The agent's structured decision logs provide a natural foundation for future fine-tuning efforts, and the modular architecture supports incremental extension to cover the full game.

% --------------------------------------------------------------------------- %

\bibliographystyle{plainnat}
\bibliography{references}

\end{document}
